%! Author = lili
%! Date = 6/22/26

\section{Tutorium 22.06.2026}\label{ch:tutorium-22.06.2026}
\untit{Präsenzaufgaben 2.3: Grenzwertsätze II \hfill 22.--26.\ Juni 2026}
\txo{Bearbeiten Sie auf jeden Fall Aufgabe 2.3.4. Falls Ihnen Zeit bleibt, bearbeiten Sie
zusätzlich die weiteren Aufgaben.}
\subsection{PÜ 2.3.4 \teils}\label{sec:pu-2.3.4}
% TODO Bilder Max
Bei einer Maschine, die Elektronikbauteile produziert, ist im Schnitt jedes 50. Bauteil fehlerhaft.
Wir testen eine Lieferung von 300 Bauteilen.
Wie groß ist die Wahrscheinlichkeit, dass genau 3 Bauteile fehlerhaft sind?

Denken Sie daran, einen Wahrscheinlichkeitsraum anzugeben und die benötigten Zufallsvariablen
zu definieren!

\begin{auf}
    \[
        \eraum =\{0,1\}^n, \quad n = 300
    \]
    mit
    \[
        \wfk(\ele) = p^{\sum_{i=1}^k \ele_i} q^{n - \sum_{i=1}^k \ele_i} \quad mit \quad  p = \frac{1}{50}, q =  \frac{49}{50}
    \]
    Teil ist kaputt ist Erfolg, damit wir uns nicht mit Gegenwahrscheinlichkeiten auseinandersetzen müssen.
    \[
        S_n(\ele) = \sum_{i=1}^n \ele_i
    \]
    Aunahme: teile sind Unabhängig voneinander kaputt und immer $p = \frac{1}{50}$
    \begin{align*}
        \pe(S_n = x) & = \binom{n}{x} \cdot p^3 \cdot q^{n-3} \\
        & = \binom{300}{3} (\frac{1}{50})^3 (\frac{49}{50})^{297} \\
        & \approx 0,0883
    \end{align*}
    Poisson Approximation (p klein q groß)
    \[\alpha = np = \frac{300}{50} = 6\]
    Sei $X \sim Poi(\alpha)$
    \begin{align*}
        \pe(S_n = 3) & \approx \pe(X=3) \\
        & = \frac{\alpha^3}{3!} e^{-\alpha} \\
        & = \frac{6^3}{3!} e^{-6} \\
        & \approx 0,0892
    \end{align*}
    \begin{align*}
        \pe(S_n = k ) & = \binom{n}{k} p^k q^{n-k} \\
        \pe(S_n \in A) & = \sum_{k\in A} \pe(S_n = k) \\
        & = \sum_{k\in A}
    \end{align*}
    % TODO
    TODO ! Noch Rest von Bildern eintragen
\end{auf}

\subsection{PÜ 2.3.5 \nochnichts (Lösung da, noch nicht eingetragen)}
% TODO Bilder Max
%Ein Buch habe $1000$ Seiten und bestehe aus $50$ Kapiteln zu je $20$ Seiten. Außerdem
%gibt es in dem Buch $40$ Druckfehler. Wir nehmen an, dass die Druckfehler unabhängig und
%mit gleicher Wahrscheinlichkeit auf den Seiten auftreten. Die Wahrscheinlichkeit, dass
%ein Druckfehler auf einer bestimmten Seite auftritt, sei also für jeden Druckfehler
%$\tfrac{1}{1000}$. Um die Druckfehler zu beseitigen, wird das Buch von einem Lektor
%geprüft. Leider findet dieser keinen einzigen Fehler, sondern erzeugt weitere
%Druckfehler. In den ersten $10$ Kapiteln je Kapitel einen, in den Kapiteln $11$ bis $40$
%jeweils $3$ und in den Kapiteln $41$ bis $50$ sogar $5$ Fehler je Kapitel. Wir wählen nun
%zufällig ein Kapitel.
%
%\begin{auf}[Geben Sie einen geeigneten Wahrscheinlichkeitsraum an, um die folgenden
%Aufgabenteile bearbeiten zu können. Diskutieren Sie insbesondere, welche
%Annahmen Sie treffen müssen.]
%    NOCH KEINE LÖSUNG % TODO NOCH KEINE LÖSUNG
%\end{auf}
%
%\begin{auf}[Definieren Sie eine Zufallsvariable $F$, die angibt, wie viele Druckfehler
%insgesamt in dem ausgewählten Kapitel auftreten.]
%    NOCH KEINE LÖSUNG % TODO NOCH KEINE LÖSUNG
%\end{auf}
%
%\begin{auf}[Berechnen Sie die Wahrscheinlichkeit, dass in dem zufällig gewählten Kapitel
%$10$ Druckfehler auftauchen. Verwenden Sie dazu eine geeignete Approximation.]
%    NOCH KEINE LÖSUNG % TODO NOCH KEINE LÖSUNG
%\end{auf}

\subsection{PÜ 2.3.6 \nochnichts (Lösung da, noch nicht eingetragen)}
% TODO Bilder Max
%Gerd und Andrea gehen an denselben $200$ Arbeitstagen pro Jahr in ihrer jeweiligen
%Mittagspause unabhängig voneinander in das zentral gelegene Café Spitzendeckchen.
%Sowohl Gerd als auch Andrea treffen mit uniformer Verteilung zwischen 12:00 und 13:00 Uhr
%ein und verbringen jeweils genau eine Viertelstunde im Café, bevor sie wieder gehen.
%
%\begin{auf}[Bestimmen Sie die Wahrscheinlichkeit, dass die beiden sich an einem bestimmten
%Tag (z.B.\ am ersten Arbeitstag des Jahres) treffen.]
%    NOCH KEINE LÖSUNG % TODO NOCH KEINE LÖSUNG
%\end{auf}
%
%\begin{auf}[Berechnen Sie die erwartete Anzahl an Treffen im Jahr.]
%    NOCH KEINE LÖSUNG % TODO NOCH KEINE LÖSUNG
%\end{auf}
%
%\begin{auf}[Approximieren Sie mit Hilfe des Satzes von de Moivre--Laplace die
%Wahrscheinlichkeit, dass es zu zwischen $80$ und $90$ Treffen kommt (beide
%Grenzen eingeschlossen).]
%    NOCH KEINE LÖSUNG % TODO NOCH KEINE LÖSUNG
%\end{auf}
%
%\begin{auf}[Approximieren Sie mit Hilfe des Satzes von de Moivre--Laplace die
%Wahrscheinlichkeit, dass es zu mindestens $90$ Treffen kommt (beide Grenzen
%eingeschlossen).]
%    NOCH KEINE LÖSUNG % TODO NOCH KEINE LÖSUNG
%\end{auf}
%
%\txo{Denken Sie daran, einen Wahrscheinlichkeitsraum anzugeben und die benötigten
%Zufallsvariablen zu definieren!}